
\chapter{Introduction} 
\label{ch:Chapter1}
\vfill \minitoc \newpage

Este documento consolida o fecho do escopo do módulo atual do Nature Protector, as
decisões arquiteturais relevantes, a decomposição funcional do módulo, o modelo es-
trutural mínimo e os respetivos critérios de aceitação. A sua finalidade é servir de base
intermédia entre a proposta de projeto e o relatório final, permitindo alinhar implementa-
ção, integração, demonstração e validação sem reabrir decisões estruturais que já devem
ficar estabilizadas no final da Fase 0.

O documento não substitui o relatório final. Funciona, antes, como especificação
técnica de referência do módulo de prevenção, com foco naquilo que é necessário para
construir, integrar, demonstrar e avaliar o sistema. Por esse motivo, privilegia clareza
arquitetural, controlo de âmbito, rastreabilidade, coerência entre decisões e utilidade
prática para o desenvolvimento.